% ===========================================================================
% SECTION 6 CLOSED-LOOP RESULTS
% Target: 1200-1500 words
% ===========================================================================

\section{Closed-Loop Results}
\label{sec:closed_loop_results}

This section reports two comparisons under two different protocols,
and this ordering matters for what can and cannot be concluded from
each. Section~\ref{subsec:unified_protocol} presents the primary,
cross-architecture comparison: a single unified protocol across all
seven controllers, resolving several inconsistencies identified in an
earlier draft (differing simulation windows, three different
statements of the prediction/control horizon, and
\texttt{predict()}-based state functions for some architectures but
not others -- the last of which was found to silently corrupt solver
feasibility via single-precision gradient truncation, per
Section~\ref{sec:mpc_framework}). Sections~\ref{subsec:v1_results}
and~\ref{subsec:v2_results} retain the earlier, separately-run V1 and
V2 snapshots for reference and for within-snapshot comparisons that
remain valid (e.g., PINN vs.\ Baseline within V1), but these two
snapshots use different windows, horizons, and (for V1's GP) a
different active-set size than the unified protocol, and should not
be used for cross-architecture claims spanning both snapshots.

\subsection{Unified-Protocol Comparison (All Architectures)}
\label{subsec:unified_protocol}

\hspace{1.5cm} Table~\ref{tab:unified_protocol} reports all seven controllers under a
single protocol: $T=\SI{60}{s}$, $N_p=15$, $N_c=4$, $T_s=\SI{0.1}{s}$,
$V=\SI{14}{m/s}$, seed $2025$, with every state function computed via
the verified manual (double-precision) forward pass rather than
\texttt{predict()} throughout. GP-v2's uncertainty term is included as
an explicit ablation ($\kappa=0$, the standard cost with no
uncertainty penalty, versus $\kappa=0.8$, the full cost) against an
otherwise identical base cost, rather than reported as a single
entangled configuration. LSTM is excluded from this table on the same
grounds as in the V1 snapshot below: its \texttt{predict()} cost is
already precisely quantified in Table~\ref{tab:predict_bypass} and a
full 600-step run adds no new information at a cost of several hours.

\begin{table}[H]
\centering
\caption{Closed-loop results under the unified protocol (all
controllers, manual forward pass, GP-v2 uncertainty ablation).}
\label{tab:unified_protocol}
\small
\begin{tabular}{p{3cm}ccccc}
\toprule
Controller & RMSE$_\omega$ (rpm) & Pitch act.\ (deg) & Mean $C_p$ & CPU (ms/step) & Infeasible \\
\midrule
Gain-Sched.\ PI (industrial ref.) & 0.613 & 71.6 & 0.2560 & 0.05 & N/A \\
Baseline (perfect-model MPC) & 0.671 & 272.7 & 0.2365 & 23.7 & 0/600 \\
LLNFM & 0.602 & 477.8 & 0.2503 & 3282.1 & 0/600 \\
GP-v2 ($\kappa=0$) & 0.604 & 0.0 & 0.2498 & 8604.5 & 600/600 \\
GP-v2 ($\kappa=0.8$) & 0.604 & 0.0 & 0.2498 & 8025.2 & 600/600 \\
PINN-v2 & 0.966 & 158.7 & 0.3380 & 21.4 & 0/600 \\
TCN & 0.717 & 322.9 & 0.2775 & 106.0 & 0/600 \\
SW-MLP & 0.757 & 401.0 & 0.2264 & 55.2 & 0/600 \\
\bottomrule
\end{tabular}
\end{table}

\hspace{1.5cm} A gain-scheduled PI pitch controller, reproducing the NREL 5MW
reference turbine's published baseline controller
\cite{jonkman2009}, is included as the genuine industrial reference
this paper's introduction motivates MPC against (item 2.7 of the
reproducibility review). We use the published gains directly
($K_{P,0}=0.01882681$~s, $K_{I,0}=0.008068634$, gain-correction corner
angle $6.30\deg$) rather than re-tuning them, since this plant's own
parameters (rotor radius, gearbox ratio, air density, rated generator
torque, and rated speed) match the published NREL 5MW reference
closely enough for this to be meaningful -- notably, the published
rated generator speed divided by this plant's gearbox ratio
($122.910/97=\SI{1.2671}{rad/s}$) reproduces this project's own
$\omega_r$ exactly.\\ \textbf{This label change matters}: the
controller previously called simply ``Baseline'' throughout this
paper is predictive control with the \emph{true, undiscretized plant
model} substituted for a learned surrogate -- an upper bound on
achievable MPC performance, not the industrial reference an earlier
draft's introduction implied it represented. It is relabeled
``Baseline (perfect-model MPC)'' throughout this paper accordingly.
The gain-scheduled PI controller achieves \emph{both} lower tracking
RMSE ($0.613$ vs.\ $\SI{0.671}{rpm}$) \emph{and} dramatically lower
pitch activity ($71.6$ vs.\ $272.7\deg$) than the perfect-model MPC
baseline, at negligible computational cost
($\SI{0.05}{ms}$/step). This is a meaningful finding in its own
right: it indicates that the perfect-model MPC baseline's own tuning
(cost weights, horizon) leaves performance on the table relative to a
well-tuned classical reference, and that any surrogate-based MPC
controller in this study should be judged against \emph{both}
references, not the perfect-model MPC baseline alone. Rotor speed
under the gain-scheduled PI ranges $[10.70, 13.31]$~rpm -- comfortably
within normal above-rated operation, and substantially better
regulated than the RMSE alone might suggest.

\hspace{1.5cm} The $\kappa=0$ and $\kappa=0.8$ rows are numerically identical: GP-v2's
total closed-loop infeasibility ($600/600$, frozen actuator, see
Section~\ref{sec:mpc_framework}) persists whether or not the
uncertainty penalty is present at all, confirming that the failure is
structural to \texttt{cost\_gp.m}'s formulation rather than an
artifact of the uncertainty term's specific weight. PINN-v2's figures
here ($0.966$~rpm, $158.7\deg$) closely reproduce the independently
generated $N_c=4$ sensitivity-sweep result ($0.965$~rpm, $157.0\deg$,
Section~\ref{sec:mpc_framework}), a useful cross-check that this table
is internally consistent. Every controller except GP-v2 achieves
$0/600$ infeasible steps under this protocol. A caveat applies to
every number in this table, however: all controllers use
\texttt{MaxIterations}$=30$, under which \texttt{ExitFlag} is
formally non-positive at nearly every step for every controller
(Table~\ref{tab:maxiter_sensitivity}, Section~\ref{sec:mpc_framework}),
meaning the solver stops short of its own optimality tolerance
throughout. A direct sensitivity sweep confirms this does not
materially affect the results reported here: RMSE is stable from
\texttt{MaxIterations}$=30$ onward for the Baseline and two
representative surrogates tested, with only
\texttt{MaxIterations}$=10$ producing a materially different, worse
outcome.


\hspace{1.5cm} Earlier, separately-run V1 and V2 snapshots (different windows,
horizons, and, for GP, a different active-set size than the unified
protocol above) are retained in Appendix~\ref{sec:appendix_legacy_runs}
for reference and reproducibility, since they document this paper's
own diagnostic history (including the since-resolved TCN/PINN-v2
numerical-coincidence finding discussed in
Section~\ref{sec:mpc_framework}), but are no longer this paper's
primary reported result -- superseded by the single unified protocol
of Section~\ref{subsec:unified_protocol} above (item 5.7 of the
reproducibility review).